\section{Methodik und Vorgehen}
\label{sec:methodik}

\subsection{Forschungsmethodischer Ansatz}
\label{sec:methodik-intro}
\hl{stucture start}
\begin{itemize}
    \item Annahmen aufzeigen auf welchen die Arbeit aufbaut
    \item Strukturierung des Vorgehen gem Drawio Prozess → Grundlage basierend auf Enterprise Architektur-Buch. Quasi von Business hinzu Anforderungen kommen, deshalb haben wir das mittel der capability map gewählt etc.
\end{itemize}
\hl{structure end}

\subsection{Vorgehen zur Identifikation von Landingzone Komponenten}
\label{sec:methodik-landingzones}
\hl{stucture start}
\begin{itemize}
    \item Kriterien beschreiben mit welchen die verschiedenen Provider analysiert wurden, bspw Aspekte und so
    \item Detaillierte Beschreibung der Erkenntnis bzw Tabelle
\end{itemize}
\hl{structure end}

\subsection{Synthese eines abstrahierten Metamodells}
\label{sec:methodik-metamodell}
\hl{stucture start}
\begin{itemize}
    \item Metamodell definieren basierend auf Fachbegriffen und Domänenbegriffen
\end{itemize}
\hl{structure end}

\subsection{Erhebung von Landing-Zone Fähigkeiten}
\label{sec:methodik-capability}
\hl{stucture start}
\begin{itemize}
    \item Quasi als Folge → Erstellen Capability Map inkl Auswahl durch den Auftraggeber (Iteration mit L1/L2 Strukturierung)
    \item Priorisierung der Capabilities mit Auftraggeber
\end{itemize}
\hl{structure end}

\subsection{Strukturierung der Evaluation der Architektur}
\label{sec:methodik-evaluation}
\hl{stucture start}
\begin{itemize}
    \item Evaluationskriterien für Traceability Architektur definieren
    \item Use-Cases definieren / Ausprägung
\end{itemize}
\hl{structure end}

